\documentclass[11pt,a4paper]{article}

\usepackage[T1]{fontenc}
\usepackage[utf8]{inputenc}
\usepackage{lmodern}
\usepackage{amssymb,amsmath,amsthm,mathtools}
\usepackage[margin=2.5cm]{geometry}
\usepackage{parskip}
\usepackage[hidelinks]{hyperref}
\usepackage{microtype}

% Theorem environments (arXiv / journal style)
\theoremstyle{plain}
\newtheorem{theorem}{Theorem}[section]
\newtheorem{proposition}[theorem]{Proposition}
\newtheorem{lemma}[theorem]{Lemma}
\newtheorem{corollary}[theorem]{Corollary}

\theoremstyle{definition}
\newtheorem{definition}[theorem]{Definition}
\newtheorem{example}[theorem]{Example}
\newtheorem{remark}[theorem]{Remark}

% Notation
\newcommand{\Shape}{\mathrm{Shape}}
\newcommand{\Pos}{\mathrm{Pos}}
\newcommand{\Type}{\mathbf{Type}}
\newcommand{\id}{\mathrm{id}}
\newcommand{\src}{\mathrm{src}}
\newcommand{\tgt}{\mathrm{tgt}}
\newcommand{\rootpos}{\mathrm{root}}
\newcommand{\sub}{\mathrm{sub}}
\newcommand{\shift}{\mathrm{shift}}

\title{Shipping Networks as Directed Containers}
\author{Extracted from a machine-checked development}
\date{}

\begin{document}
\maketitle

\begin{abstract}
Every directed (multi)graph~$G$ determines a directed container whose shapes
are the nodes of~$G$ and whose positions at a node~$v$ are the finite directed
paths starting at~$v$, recorded together with their endpoints. The root is the
empty path, the sub-shape map returns the endpoint of a path, and the shift is
path concatenation. As a category this construction is the free category
on~$G$. We record the elementary path algebra, the directed-container laws, and
a concrete running example---an eleven-facility wine shipping network---with
explicit route counts. The proofs of all stated results have been
machine-checked; we record only the mathematical content.
\end{abstract}

\tableofcontents
\bigskip

%======================================================================
\section{Introduction}
%======================================================================

In the theory of directed containers, a structure of shapes, positions, roots,
and shifts axiomatises a free category on a set of generators. A natural source
of examples is directed graphs: nodes generate shapes, and directed paths
generate positions. The present note isolates that construction.

Given a directed multigraph~$G$ (allowing parallel edges and loops, with no
finiteness assumption), one forms a directed container $\mathrm{network}(G)$
by taking shapes to be nodes and positions at~$v$ to be pairs $\langle w,\rho
\rangle$ consisting of an endpoint~$w$ and a finite directed path $\rho$ from
$v$ to~$w$. Concatenation of paths supplies the monoidal structure on positions
required by the directed-container axioms. Viewed as a category (via the free
category associated to a directed container), $\mathrm{network}(G)$ is the free
category on~$G$.

The development works throughout in a single universe of types. The statements
below are those of a formal development in dependent type theory; proofs are
indicated by the phrase ``by lean,'' meaning they have been checked by a proof
assistant rather than written out in prose.

\paragraph{What this note claims.}
In ordinary mathematical language: the free category on a directed graph admits
a presentation as a directed container whose positions are path-endpoint pairs.
The directed-container laws reduce almost definitionally to the unit and
associativity laws for path concatenation. On a concrete acyclic shipping
graph with eleven facilities, the construction is inhabited and yields the
expected route cardinalities (nine routes out of the first vineyard; a single
trivial route at each retail sink).

%======================================================================
\section{Directed graphs and paths}
%======================================================================

\begin{definition}[Directed graph]\label{def:digraph}
A \textbf{directed (multi)graph} $G$ consists of
\begin{enumerate}
  \item a type $V$ of \emph{nodes} (vertices),
  \item a type $E$ of \emph{edges}, and
  \item endpoint maps $\src : E \to V$ and $\tgt : E \to V$.
\end{enumerate}
Parallel edges and loops are allowed. No finiteness of $V$ or $E$ is assumed.
\end{definition}

\begin{definition}[Directed path]\label{def:path}
Let $G$ be a directed graph. For nodes $v,w \in V$, the type $\mathrm{Path}_G(v,w)$
of \textbf{finite directed paths} from $v$ to $w$ is defined inductively by
\begin{enumerate}
  \item $\mathrm{nil}(v) : \mathrm{Path}_G(v,v)$ (the empty path at~$v$), and
  \item if $e \in E$ and $p : \mathrm{Path}_G\bigl(\tgt(e),\, w\bigr)$, then
        \[
          \mathrm{cons}(e,p) \;:\; \mathrm{Path}_G\bigl(\src(e),\, w\bigr)
        \]
        (follow $e$, then continue along~$p$).
\end{enumerate}
\end{definition}

Thus paths are ``cons-at-the-front'': an outgoing edge followed by a path from
its target.

\begin{definition}[Path concatenation]\label{def:append}
Given paths $p : \mathrm{Path}_G(v,w)$ and $q : \mathrm{Path}_G(w,u)$, their
\textbf{concatenation} $p \mathbin{{+}{+}} q : \mathrm{Path}_G(v,u)$ is defined
by recursion on~$p$:
\begin{align}
  \mathrm{nil}(v) \mathbin{{+}{+}} q &\;=\; q, \\
  \mathrm{cons}(e,p) \mathbin{{+}{+}} q &\;=\; \mathrm{cons}\bigl(e,\, p \mathbin{{+}{+}} q\bigr).
\end{align}
\end{definition}

\begin{proposition}[Left unit]\label{prop:nil-append}
For every path $p : \mathrm{Path}_G(v,w)$,
\[
  \mathrm{nil}(v) \mathbin{{+}{+}} p \;=\; p.
\]
\end{proposition}

\begin{proof}
By lean.
\end{proof}

\begin{proposition}[Right unit]\label{prop:append-nil}
For every path $p : \mathrm{Path}_G(v,w)$,
\[
  p \mathbin{{+}{+}} \mathrm{nil}(w) \;=\; p.
\]
\end{proposition}

\begin{proof}
By lean.
\end{proof}

\begin{proposition}[Associativity]\label{prop:append-assoc}
For paths $p : \mathrm{Path}_G(v,w)$, $q : \mathrm{Path}_G(w,u)$, and
$r : \mathrm{Path}_G(u,x)$,
\[
  \bigl(p \mathbin{{+}{+}} q\bigr) \mathbin{{+}{+}} r
  \;=\;
  p \mathbin{{+}{+}} \bigl(q \mathbin{{+}{+}} r\bigr).
\]
\end{proposition}

\begin{proof}
By lean.
\end{proof}

The left unit law holds definitionally. The right unit and associativity laws
are proved by induction on the first path.

%======================================================================
\section{The network directed container}
%======================================================================

We recall the data of a directed container, specialised to the form used below.
A directed container consists of a type of shapes, a family of position types,
a distinguished root position at each shape, a sub-shape map sending a shape
and a position to a shape, and a shift operation composing positions, subject
to five equational laws (identity and associativity for shift, together with
compatibility of sub with root and shift).

\begin{definition}[Network directed container]\label{def:network}
Let $G$ be a directed graph. The \textbf{network directed container}
$\mathrm{network}(G)$ is defined by
\begin{align}
  \Shape\bigl(\mathrm{network}(G)\bigr)
  &\;=\; V, \\[4pt]
  \Pos\bigl(\mathrm{network}(G)\bigr)(v)
  &\;=\;
  \sum_{w \,:\, V}\; \mathrm{Path}_G(v,w), \\[4pt]
  \rootpos(v)
  &\;=\;
  \bigl\langle v,\, \mathrm{nil}(v) \bigr\rangle, \\[4pt]
  \sub\bigl(v,\, \langle w,\rho \rangle\bigr)
  &\;=\;
  w, \\[4pt]
  \shift\bigl(v,\, \langle w,\rho \rangle,\, \langle u,\sigma \rangle\bigr)
  &\;=\;
  \bigl\langle u,\, \rho \mathbin{{+}{+}} \sigma \bigr\rangle.
\end{align}
Here a position at~$v$ is a dependent pair $\langle w,\rho \rangle$ of an
endpoint and a path from~$v$ to that endpoint; the sub-shape is the first
projection; the root is the empty path; and shift is path concatenation.
\end{definition}

\begin{remark}[Directed-container laws]
With this packaging, several of the directed-container axioms hold by
definitional equality: the endpoint of the empty path is the node itself, and
the endpoint of a concatenation is the endpoint of the second path. The
remaining laws are precisely the right unit law for concatenation
(Proposition~\ref{prop:append-nil}) and associativity
(Proposition~\ref{prop:append-assoc}), each transported along the dependent
pair of endpoint and path. Indexing paths by both endpoints---rather than by
start alone---is what makes the sub-shape a pure first projection and collapses
the transport-heavy equalities to the path algebra above.
\end{remark}

\begin{proposition}[Network is a directed container]\label{prop:network-dc}
For every directed graph~$G$, the data of Definition~\ref{def:network}
satisfy the directed-container axioms.
\end{proposition}

\begin{proof}
By lean.
\end{proof}

\begin{remark}
Finiteness of~$G$ is not assumed. Cyclic graphs yield infinitely many
positions at a node; that phenomenon is treated separately in the graph-rung
analysis of the associated development.
\end{remark}

%======================================================================
\section{Running example: the Sawasdee Cellars wine network}
%======================================================================

We instantiate the construction on a concrete shipping graph: three South
Australian vineyards ship through an export agency and a Melbourne--Chonburi
sea-freight leg to two Thai distributors and three retail shops. The graph is
a polytree (thinness is treated elsewhere); here we only form the network and
count routes.

\begin{definition}[Sawasdee facilities and legs]\label{def:wine}
Let $\mathrm{Facility}$ be the eleven-element type
\begin{align*}
  &\mathrm{vineyard}_1,\;
  \mathrm{vineyard}_2,\;
  \mathrm{vineyard}_3,\;
  \mathrm{sawasdee},\;
  \mathrm{melbournePort},\;
  \mathrm{thaiPort},\\
  &\mathrm{distributor}_1,\;
  \mathrm{distributor}_2,\;
  \mathrm{retail}_1,\;
  \mathrm{retail}_2,\;
  \mathrm{retail}_3,
\end{align*}
and let $\mathrm{Leg}$ be the ten-element type of shipping legs
\begin{align*}
  &v_1s,\; v_2s,\; v_3s
    &&\text{(vineyards $\to$ Sawasdee)}, \\
  &smp
    &&\text{(Sawasdee $\to$ Melbourne Port)}, \\
  &\mathrm{sea}
    &&\text{(Melbourne Port $\to$ Thai Port)}, \\
  &td_1,\; td_2
    &&\text{(Thai Port $\to$ distributors)}, \\
  &d_1r_1,\; d_1r_2
    &&\text{(Distributor$_1$ $\to$ Retail$_1$, Retail$_2$)}, \\
  &d_2r_3
    &&\text{(Distributor$_2$ $\to$ Retail$_3$)}.
\end{align*}
Write $G_{\mathrm{wine}}$ for the directed graph with these nodes and edges and
the evident source and target maps, and write
$\mathcal{N}_{\mathrm{wine}} := \mathrm{network}(G_{\mathrm{wine}})$ for the
associated directed container.
\end{definition}

A \textbf{route} out of a facility~$v$ is a position of
$\mathcal{N}_{\mathrm{wine}}$ at~$v$, i.e.\ a pair $\langle w,\rho \rangle$
with $\rho : \mathrm{Path}_{G_{\mathrm{wine}}}(v,w)$. Given a leg~$e$ and a
route out of $\tgt(e)$, one obtains a route out of $\src(e)$ by consing~$e$
onto the underlying path.

\begin{proposition}[Route counts]\label{prop:route-counts}
In $G_{\mathrm{wine}}$:
\begin{enumerate}
  \item There are exactly nine routes out of $\mathrm{vineyard}_1$
        (the empty route, together with one route to each of the eight
        facilities reachable from $\mathrm{vineyard}_1$).
  \item There is exactly one route out of $\mathrm{retail}_1$ (the empty
        route: a retail shop is a sink).
  \item The endpoints of the nine routes out of $\mathrm{vineyard}_1$ are,
        in order of enumeration,
        \[
          \begin{aligned}
            &\mathrm{vineyard}_1,\;
            \mathrm{sawasdee},\;
            \mathrm{melbournePort},\;
            \mathrm{thaiPort},\\
            &\mathrm{distributor}_1,\;
            \mathrm{retail}_1,\;
            \mathrm{retail}_2,\;
            \mathrm{distributor}_2,\;
            \mathrm{retail}_3.
          \end{aligned}
        \]
\end{enumerate}
\end{proposition}

\begin{proof}
By lean.
\end{proof}

\begin{remark}
The enumerations are constructed in reverse topological order, so that each
facility's route list is assembled from the lists of facilities strictly
downstream. Completeness of the lists for this acyclic graph is immediate by
inspection; the formal checks assert only the stated lengths and endpoints.
\end{remark}

\begin{proposition}[Root and sub on a concrete route]\label{prop:full-route}
Let $\rho$ be the full five-leg route
$\mathrm{vineyard}_1 \to \mathrm{sawasdee} \to \mathrm{melbournePort} \to
\mathrm{thaiPort} \to \mathrm{distributor}_1 \to \mathrm{retail}_1$. Then:
\begin{enumerate}
  \item the empty route at $\mathrm{vineyard}_1$ coincides with the root
        position of $\mathcal{N}_{\mathrm{wine}}$ at $\mathrm{vineyard}_1$;
  \item $\sub(\mathrm{vineyard}_1,\rho) = \mathrm{retail}_1$;
  \item shifting the root by~$\rho$ recovers~$\rho$.
\end{enumerate}
\end{proposition}

\begin{proof}
By lean.
\end{proof}

Thus on the running example, $\sub$ is literally ``where the route ends,'' and
shift specialises to path concatenation of an empty prefix with a concrete
route.

%======================================================================
\section{Discussion}
%======================================================================

\paragraph{What has been established.}
Every directed graph presents a directed container whose positions are
path-endpoint pairs and whose monoidal structure is concatenation. The
construction is free on the graph: as a category it is the free category
generated by the edges. The packaging with double-indexed paths makes most of
the directed-container equalities definitional, leaving only the classical unit
and associativity laws of path composition.

\paragraph{Scope.}
The development does not assume finiteness, does not prove thinness of polytree
networks (treated separately), and does not address trackings, aggregations, or
other structure on top of $\mathrm{network}(G)$. Those refinements build on the
same foundations.

\paragraph{Machine checking.}
Every definition and theorem statement above corresponds to a declaration in a
machine-checked formalisation. The phrase ``by lean'' indicates that the proof
has been verified by a proof assistant; we have not reproduced the formal
proof scripts here.

%======================================================================
\section*{Acknowledgements}
%======================================================================

The mathematical content follows the construction of shipping networks as
directed containers in the supply-chain development that motivated this note.
The presentation is a prose rendering of a formal development intended for
human review of the theorem statements.

\begin{thebibliography}{9}

\bibitem{abbott2005}
M.~Abbott, T.~Altenkirch, and N.~Ghani.
\newblock Containers: constructing strictly positive types.
\newblock \emph{Theoretical Computer Science}, 342(1):3--27, 2005.

\bibitem{ahman2014}
D.~Ahman and T.~Uustalu.
\newblock Coalgebraic update lenses.
\newblock In \emph{MFPS}, volume~308 of \emph{ENTCS}, pages 25--48, 2014.

\end{thebibliography}

\end{document}
